\documentclass[openany, fontset=none]{ctexbook}
\ctexset{
  chapter={
    beforeskip=0pt,
    afterskip=20pt,
   }
}
\setmainfont{XITS}
% 中文默认字体：方正书宋_GBK，粗体为思源宋体半粗体，斜体为方正楷体_GBK
\setCJKmainfont{FZNewShuSong-Z10}[BoldFont={Source Han Serif SC SemiBold}, ItalicFont=FZKai-Z03]
% 中文无衬线字体：方正黑体_GBK，粗体为思源黑体中粗体
\setCJKsansfont{FZHei-B01}[BoldFont={Source Han Sans CN Medium}]
% 中文等宽字体：方正仿宋_GBK
\setCJKmonofont{FZFangSong-Z02}

\newCJKfontfamily\songti{FZNewShuSong-Z10}[BoldFont={Source Han Serif SC SemiBold}]
\newCJKfontfamily\xbsong{Source Han Serif SC SemiBold} % 小标宋
\newCJKfontfamily\dbsong{Source Han Serif SC Bold} % 大标宋
\newCJKfontfamily\cusong{Source Han Serif SC Heavy} % 粗宋
\newCJKfontfamily\heiti{FZHei-B01}[BoldFont={Source Han Sans CN Medium}]
\newCJKfontfamily\dahei{Source Han Sans CN Medium} % 大黑
\newCJKfontfamily\cuhei{Source Han Sans CN Bold} % 粗黑
\newCJKfontfamily\fangsong{FZFangSong-Z02}
\newCJKfontfamily\kaiti{FZKai-Z03}
\usepackage{layout}
\usepackage[a4paper, twoside=false, includemp]{geometry}
\geometry{left=1cm, right=0.1cm, bottom=1.6cm, top=1.6cm, marginparwidth=36pt}
\usepackage[notindex]{tocbibind}
\usepackage{xcolor}
\definecolor{mygreen}{RGB}{3,143,94}
\colorlet{bgcolor}{mygreen!5}
\usepackage{zhlineskip}
\usepackage{multirow}
\usepackage{makecell}
\usepackage{graphicx}
\usepackage{amsmath}
\usepackage[edges]{forest}
\forestset{
  default preamble={tikz+={\tikzset{
    show background rectangle,
    background rectangle/.style={draw=mygreen}
  }}},
  my tree/.style={
      for tree={
          grow'=east,
          forked edge,
          edge={line, mygreen},
          thick,
          draw=mygreen,
          fill=bgcolor,
          rounded corners=4pt,
          font=\small,
        }
    }
}
\usepackage{tikz}
% quiver
\usepackage{tikz-cd}
\usepackage{amssymb}
\usetikzlibrary{arrows.meta, backgrounds}
\usetikzlibrary{calc}
\usetikzlibrary{decorations.pathmorphing}
\tikzset{line/.style={draw, -Stealth}}
% \tikzset{line/.style={draw, -{Circle[open]}}}
\tikzset{
  curve/.style={
      settings={#1},
      to path={
          (\tikztostart) .. controls ($(\tikztostart)!\pv{pos}!(\tikztotarget)!\pv{height}!270:(\tikztotarget)$)
          and ($(\tikztostart)!1-\pv{pos}!(\tikztotarget)!\pv{height}!270:(\tikztotarget)$)
          .. (\tikztotarget)\tikztonodes
        },
    },
  settings/.code={
      \tikzset{quiver/.cd, #1}
      \def\pv##1{\pgfkeysvalueof{/tikz/quiver/##1}}
    },
  quiver/.cd, pos/.initial=0.35, height/.initial=0,
}
\tikzset{tail reversed/.code={\pgfsetarrowsstart{tikzcd to}}}
\tikzset{2tail/.code={\pgfsetarrowsstart{Implies[reversed]}}}
\tikzset{2tail reversed/.code={\pgfsetarrowsstart{Implies}}}
\tikzset{no body/.style={/tikz/dash pattern=on 0 off 1mm}}
\tikzcdset{
  background color=bgcolor,
  every label/.append style={font=\small},
  marking/.append style={font=\small},
  description/.append style={font=\small},
}
\tikzset{
  mytest/.style={
      show background rectangle,
      background rectangle/.style={
          draw=mygreen,
          fill=bgcolor,
          rounded corners=4pt,
        },
      font=\small,
    },
}
% end quiver
\usepackage[all, pdf]{xy}
\usepackage{imakeidx}
\makeindex[
  intoc=true,
  columns=1
]
\usepackage{nameref}
\usepackage[bookmarksopen=true]{hyperref}
\usepackage{bookmark}
\usepackage{paracol}
\usepackage{tcolorbox}
\tcbuselibrary{skins}
\tcbuselibrary{breakable}
\newtcolorbox{kk}[2][]{
  % padding left/right
  left=4pt,
  right=4pt,
  % margin top/botom
  before skip=6pt,
  after skip=8pt,
  % 圆角
  arc=4pt,
  sharp corners=northwest,
  boxrule=.5pt, % border
  colframe=mygreen, % 线框颜色
  colback=mygreen!5, % 框背景色
  drop lifted shadow,
  % 标题
  enhanced, % 显示
  boxed title size=copy, % 大小
  lefttitle=4pt, % padding left
  coltitle=white, % 标题颜色
  colbacktitle=mygreen, % 标题框背景色
  % fonttitle=\bfseries,
  boxed title style={arc=0pt},
  attach boxed title to top left,
  title={#2}, #1}
\usepackage[inline]{enumitem}
\setlist{nosep}
% \setlist[enumerate, 3]{label=(\arabic*), ref=\theenumii.(\arabic*)}
% \setlist[itemize, 1]{label=$\ast$}
\usepackage{xeCJKfntef}
\newcommand\xhx[1]{\CJKunderline*{#1}} % 下划线
\newcommand\scx[1]{\CJKsout*{#1}} % 删除线

\setlength{\fboxsep}{1pt}
\newcommand\hzi[1]{\textcolor{magenta}{#1}}
\newcommand\hs[1]{\textcolor{red}{#1}}
\newcommand\zd[1]{\CJKunderline*[textformat=\color{mygreen}]{#1}} % 带下划线的重点
\newcommand\bs[1]{\CJKunderline*[textformat=\kaiti\color{red}]{#1}} % 容易忽略的知识点，背诵
\newcommand\qd[1]{\textcolor{mygreen}{#1}} % 强调
\newcommand\al[1]{\hzi{\fangsong{#1}}} % 案例 不用 texttt
\newcommand\ft[1]{\hzi{\fangsong{#1}}} % 法条

\newcommand\sy[1]{{#1}\index{#1}} % 索引
\newcommand\bq[1]{{#1}\label{#1}\index{#1}} % 标签
\newcommand\yy[1]{{#1}\ref{#1}} % 引用
\newcommand\tz[1]{\hyperref[#1]{#1}} % 跳转

\newcommand\yjt{\textbf{→} } % 右箭头
\newcommand\zyjt{↔} % 左右箭头
\newcommand\sjt{\raisebox{0.2ex}{\textbf{↑}}\,} % 上箭头
\newcommand\xjt{\textbf{↓}\,} % 下箭头
\newcommand\blh{\textasciitilde} % 波浪号
\newcommand\dy{＞ } % 大于
\newcommand\xyh{＜ } % 小于
\newcommand\ydy{≈ } % 约等于
\newcommand\bdy{≠ } % 不等于
\newcommand\dydy{≥ } % 大于等于
\newcommand\xydy{≤ } % 小于等于
\newcommand\hz{|| } % 或者
\newcommand\cuo{× } % 错
\newcommand\dui{√ } % 对

\newcommand\yi{\qd{①}.} % 圆圈 1 - 4
\newcommand\er{\qd{②}.} %
\newcommand\san{\qd{③}.} %
\newcommand\si{\qd{④}.} %
\newcommand\info[2]{
  \hypersetup{
    colorlinks,
    % bookmarksnumbered,
    pdftitle={#1},
    pdfauthor={小兑},
    pdfsubject={法考 #1},
    pdfkeywords={法考 #1}
  }
  \title{#1}
  \author{ecuplxd}
  \date{#2}
}
\makeatletter
\newcommand{\DrawLine}{%
  \begin{tikzpicture}
    \path[use as bounding box] (0,0) -- (\linewidth, 0);
    \draw[color=mygreen, dashed, dash phase=2pt]
    (0 - \kvtcb@leftlower - \kvtcb@boxsep, 0)--
    (\linewidth + \kvtcb@rightlower + \kvtcb@boxsep, 0);
  \end{tikzpicture}%
}
\makeatother
